% ============================================================
%  EXAMEN PARCIAL 3  —  Métodos Numéricos
%  Tema: Integración Numérica Definida
%  Compilar:  latexmk -pdf examen_parcial3.tex
% ============================================================
\documentclass[12pt,letterpaper]{article}

% ── Codificación e idioma ──────────────────────────────────
\usepackage[utf8]{inputenc}
\usepackage[T1]{fontenc}
\usepackage[spanish]{babel}

% ── Matemáticas ───────────────────────────────────────────
\usepackage{amsmath,amssymb,mathtools}

% ── Tablas ────────────────────────────────────────────────
\usepackage{booktabs,array,multirow}

% ── Listas ────────────────────────────────────────────────
\usepackage{enumitem}

% ── Figuras ───────────────────────────────────────────────
\usepackage{graphicx}
\usepackage{float}
\graphicspath{{./capturas/}}   % imágenes en subcarpeta capturas/

% ── Geometría ─────────────────────────────────────────────
\usepackage[top=2.2cm,bottom=2.2cm,left=2.5cm,right=2.5cm]{geometry}

% ── Colores y cajas ───────────────────────────────────────
\usepackage{xcolor,mdframed}

% ── Encabezado / pie de página ────────────────────────────
\usepackage{fancyhdr}
\pagestyle{fancy}
\fancyhf{}
\rhead{\small Métodos Numéricos}
\lhead{\small Francisco Javier Ponce Sáenz — a325000}
\cfoot{\thepage}
\renewcommand{\headrulewidth}{0.4pt}

% ── Entorno problema ──────────────────────────────────────
\newcounter{problema}
\newenvironment{problema}[1][]{%
  \refstepcounter{problema}%
  \begin{mdframed}[
      linewidth=1pt,
      linecolor=black!40,
      backgroundcolor=gray!6,
      innertopmargin=6pt,
      innerbottommargin=6pt
  ]
  \noindent\textbf{Problema \theproblema}%
  \ifx&#1&\else\ \textbf{— #1}\fi\par\smallskip
}{%
  \end{mdframed}\vspace{6pt}
}

\newenvironment{solucion}{%
  \vspace{4pt}\noindent\textbf{Solución:}\par\smallskip
}{%
  \par
}

% ── Comandos útiles ───────────────────────────────────────
\newcommand{\respuesta}[1]{\par\medskip\noindent\fbox{\textbf{Resultado:}~#1}}
\newcommand{\tol}[1]{\textrm{Tol} = #1}

% ── Colores de método ─────────────────────────────────────
\definecolor{azulmetodo}{RGB}{0,70,140}
\definecolor{verdemetodo}{RGB}{0,110,60}
\definecolor{rojometodo}{RGB}{140,20,20}
\definecolor{naranjametodo}{RGB}{180,80,0}

% ─────────────────────────────────────────────────────────
\begin{document}
% ─────────────────────────────────────────────────────────

% ══════════════════  PORTADA  ════════════════════════════
\begin{center}
  {\large\bfseries Universidad Autónoma de Chihuahua}\\[2pt]
  {\large Facultad de Ingeniería}\\[6pt]
  \rule{\linewidth}{0.8pt}\\[4pt]
  {\LARGE\bfseries Examen Parcial 3}\\[4pt]
  {\large Integración Numérica Definida}\\[4pt]
  \rule{\linewidth}{0.8pt}\\[6pt]
  {\large Métodos Numéricos}\\[4pt]
  {\large Francisco Javier Ponce Sáenz \quad|\quad a325000}\\[2pt]
  {\large\today}
\end{center}

\vspace{8pt}

% ══════════════════  FORMULARIO DE REFERENCIA  ══════════════════
\begin{mdframed}[linecolor=azulmetodo, linewidth=1.2pt, backgroundcolor=blue!4,
                  innertopmargin=8pt, innerbottommargin=8pt]
\noindent\textbf{\textcolor{azulmetodo}{Formulario — Integración Numérica}}
\vspace{4pt}

\textbf{Paso:}
\[
  h = \frac{b - a}{N}
\]

\begin{center}
\begin{tabular}{llcc}
  \toprule
  Método & Fórmula & Factor & Restricción ($N$) \\
  \midrule
  Trapecio
    & $f(a) + 2\displaystyle\sum_{i=1}^{N-1}f(x_i) + f(b)$
    & $\dfrac{h}{2}$ & cualquier $N$ \\[8pt]
  Simpson $\tfrac{1}{3}$
    & $f(a) + 4\displaystyle\sum_{\text{impar}}f(x_i) + 2\displaystyle\sum_{\text{par}}f(x_i) + f(b)$
    & $\dfrac{h}{3}$ & $N$ múltiplo de 2 \\[8pt]
  Regla $\tfrac{3}{8}$
    & $f(a) + 3\displaystyle\sum_{i\neq 3k}f(x_i) + 2\displaystyle\sum_{3k}f(x_i) + f(b)$
    & $\dfrac{3h}{8}$ & $N$ múltiplo de 3 \\[8pt]
  Boole ($\tfrac{2}{45}$)
    & $7f(a) + 32\displaystyle\sum_{\text{impar}}f(x_i) + 12\displaystyle\sum_{\substack{\text{par}\\{\not\equiv 0\,(4)}}}f(x_i) + 14\displaystyle\sum_{4k}f(x_i) + 7f(b)$
    & $\dfrac{2h}{45}$ & $N$ múltiplo de 4 \\[8pt]
  \bottomrule
\end{tabular}
\end{center}

\textbf{Richardson (Límite Diferido):}
\[
  I_{\text{LDR}} = I_h + (I_k - I_h)\cdot\frac{h^2}{h^2 - k^2}
  \qquad \bigl(k = h/2 \;\Rightarrow\; N_k = 2N_h\bigr)
\]
\end{mdframed}

\vspace{10pt}

% ══════════════════  PROBLEMA 1  ════════════════════════
\begin{problema}[Método Trapecial y Richardson]

% --- ENUNCIADO (completar con el ejercicio del examen) ---
Calcule la siguiente integral definida aplicando el \textbf{Método Trapecial} y
el \textbf{Límite Diferido de Richardson}:

\[
  \int_{a}^{b} f(x)\,dx
  \qquad\qquad
  % Sustituir a, b, f(x) y el valor de N
  a = \rule{1.5cm}{0.4pt},\quad
  b = \rule{1.5cm}{0.4pt},\quad
  N_1 = \rule{1.5cm}{0.4pt},\quad
  N_2 = \rule{1.5cm}{0.4pt}
\]

% ───── SOLUCIÓN ──────────────────────────────────────────
\begin{solucion}

\subsection*{a) Parámetros}

\[
  h_1 = \frac{b-a}{N_1} = \rule{2cm}{0.4pt},
  \qquad
  h_2 = \frac{b-a}{N_2} = \rule{2cm}{0.4pt}
\]

\subsection*{b) Tabla del Método Trapecial ($N_1$)}

\begin{center}
\begin{tabular}{c|c|c}
  \toprule
  $x_i$ & $f(x_i)$ & Coef. \\
  \midrule
  $a$  & & $\times 1$ \\[3pt]
       & & $\times 2$ \\[3pt]
       & & $\times 2$ \\[3pt]
       & & $\times 2$ \\[3pt]
       & & $\times 2$ \\[3pt]
       & & $\times 2$ \\[3pt]
  $b$  & & $\times 1$ \\[3pt]
  \midrule
  \multicolumn{2}{l|}{Sumatoria} & \rule{2cm}{0.4pt} \\
  \bottomrule
\end{tabular}
\end{center}

\[
  I_{h_1} = \frac{h_1}{2}\cdot(\text{sumatoria}) = \rule{3cm}{0.4pt}
\]

\subsection*{c) Tabla del Método Trapecial ($N_2$)}

\begin{center}
\begin{tabular}{c|c|c}
  \toprule
  $x_i$ & $f(x_i)$ & Coef. \\
  \midrule
  $a$  & & $\times 1$ \\[3pt]
       & & $\times 2$ \\[3pt]
       & & $\times 2$ \\[3pt]
       & & $\times 2$ \\[3pt]
       & & $\times 2$ \\[3pt]
       & & $\times 2$ \\[3pt]
       & & $\times 2$ \\[3pt]
       & & $\times 2$ \\[3pt]
       & & $\times 2$ \\[3pt]
       & & $\times 2$ \\[3pt]
       & & $\times 2$ \\[3pt]
       & & $\times 2$ \\[3pt]
  $b$  & & $\times 1$ \\[3pt]
  \midrule
  \multicolumn{2}{l|}{Sumatoria} & \rule{2cm}{0.4pt} \\
  \bottomrule
\end{tabular}
\end{center}

\[
  I_{h_2} = \frac{h_2}{2}\cdot(\text{sumatoria}) = \rule{3cm}{0.4pt}
\]

\subsection*{d) Límite Diferido de Richardson}

\[
  I_{\text{LDR}}
  = I_{h_1} + (I_{h_2} - I_{h_1})\cdot\frac{h_1^2}{h_1^2 - h_2^2}
  = \rule{4cm}{0.4pt}
\]

\respuesta{$I_T(N_1) = \rule{2.5cm}{0.4pt}$, \quad
           $I_T(N_2) = \rule{2.5cm}{0.4pt}$, \quad
           $I_{\text{LDR}} = \rule{2.5cm}{0.4pt}$}

\end{solucion}
\end{problema}

% ══════════════════  PROBLEMA 2  ════════════════════════
\begin{problema}[Newton-Cotes — $N = 12$]

% --- ENUNCIADO ---
Calcule la siguiente integral aplicando los \textbf{cuatro métodos de Newton-Cotes}
con $N = 12$:

\[
  \int_{a}^{b} f(x)\,dx,
  \qquad
  a = \rule{1.5cm}{0.4pt},\quad
  b = \rule{1.5cm}{0.4pt}
\]

\begin{solucion}

\subsection*{Parámetros}

\[
  N = 12, \qquad
  h = \frac{b - a}{12} = \rule{2.5cm}{0.4pt}
\]

\subsection*{Tabla de valores ($x_i$, $f(x_i)$) y coeficientes}

\begin{center}
\begin{tabular}{c|c|c|c|c|c}
  \toprule
  $i$ & $x_i$ & $f(x_i)$ & Trap. & Simp. $\tfrac{1}{3}$ & Regla $\tfrac{3}{8}$ \\
  \midrule
  0  & $a$   & & 1 & 1  & 1  \\[3pt]
  1  &       & & 2 & 4  & 3  \\[3pt]
  2  &       & & 2 & 2  & 3  \\[3pt]
  3  &       & & 2 & 4  & 2  \\[3pt]
  4  &       & & 2 & 2  & 3  \\[3pt]
  5  &       & & 2 & 4  & 3  \\[3pt]
  6  &       & & 2 & 2  & 2  \\[3pt]
  7  &       & & 2 & 4  & 3  \\[3pt]
  8  &       & & 2 & 2  & 3  \\[3pt]
  9  &       & & 2 & 4  & 2  \\[3pt]
  10 &       & & 2 & 2  & 3  \\[3pt]
  11 &       & & 2 & 4  & 3  \\[3pt]
  12 & $b$   & & 1 & 1  & 1  \\[3pt]
  \midrule
  \multicolumn{2}{c|}{Sumatoria} & &
    \rule{1.2cm}{0.4pt} & \rule{1.2cm}{0.4pt} & \rule{1.2cm}{0.4pt} \\
  \bottomrule
\end{tabular}
\end{center}

\vspace{6pt}

\begin{center}
\begin{tabular}{c|c|c}
  \toprule
  $i$ & $x_i$ & Boole ($\tfrac{2}{45}$) \\
  \midrule
  0  & $a$   & 7  \\[3pt]
  1  &       & 32 \\[3pt]
  2  &       & 12 \\[3pt]
  3  &       & 32 \\[3pt]
  4  &       & 14 \\[3pt]
  5  &       & 32 \\[3pt]
  6  &       & 12 \\[3pt]
  7  &       & 32 \\[3pt]
  8  &       & 14 \\[3pt]
  9  &       & 32 \\[3pt]
  10 &       & 12 \\[3pt]
  11 &       & 32 \\[3pt]
  12 & $b$   & 7  \\[3pt]
  \midrule
  \multicolumn{2}{c|}{Sumatoria} & \rule{1.2cm}{0.4pt} \\
  \bottomrule
\end{tabular}
\end{center}

\subsection*{Resultados por método}

\begin{center}
\begin{tabular}{lcc}
  \toprule
  Método & Fórmula aplicada & Resultado \\
  \midrule
  Trapecio
    & $I_T = \dfrac{h}{2}\cdot\Sigma$
    & \rule{3cm}{0.4pt} \\[10pt]
  Simpson $\tfrac{1}{3}$
    & $I_S = \dfrac{h}{3}\cdot\Sigma$
    & \rule{3cm}{0.4pt} \\[10pt]
  Regla $\tfrac{3}{8}$
    & $I_{3/8} = \dfrac{3h}{8}\cdot\Sigma$
    & \rule{3cm}{0.4pt} \\[10pt]
  Boole ($\tfrac{2}{45}$)
    & $I_B = \dfrac{2h}{45}\cdot\Sigma$
    & \rule{3cm}{0.4pt} \\[6pt]
  \bottomrule
\end{tabular}
\end{center}

\respuesta{Trapecio $= \rule{2cm}{0.4pt}$, \;
           Simp.\ $\frac{1}{3} = \rule{2cm}{0.4pt}$, \;
           Regla $\frac{3}{8} = \rule{2cm}{0.4pt}$, \;
           Boole $= \rule{2cm}{0.4pt}$}

\end{solucion}
\end{problema}

% ══════════════════  PROBLEMA 3 (opcional)  ════════════════════
\begin{problema}[Problema adicional]

% --- ENUNCIADO ---
\[
  \int_{a}^{b} f(x)\,dx
  \qquad
  a = \rule{1.5cm}{0.4pt},\quad
  b = \rule{1.5cm}{0.4pt},\quad
  N = \rule{1.5cm}{0.4pt}
\]

\begin{solucion}

\[
  h = \frac{b-a}{N} = \rule{2.5cm}{0.4pt}
\]

% Tabla genérica — ampliar o reducir filas según N
\begin{center}
\begin{tabular}{c|c|c|c|c|c}
  \toprule
  $i$ & $x_i$ & $f(x_i)$ & Trap. & Simp. $\tfrac{1}{3}$ & Regla $\tfrac{3}{8}$ \\
  \midrule
  & & & & & \\[3pt]
  & & & & & \\[3pt]
  & & & & & \\[3pt]
  & & & & & \\[3pt]
  & & & & & \\[3pt]
  & & & & & \\[3pt]
  & & & & & \\[3pt]
  \midrule
  \multicolumn{2}{c|}{Sumatoria} & &
    \rule{1.2cm}{0.4pt} & \rule{1.2cm}{0.4pt} & \rule{1.2cm}{0.4pt} \\
  \bottomrule
\end{tabular}
\end{center}

\respuesta{$\rule{6cm}{0.4pt}$}

\end{solucion}
\end{problema}

% ══════════════════  EVIDENCIA COMPUTACIONAL  ════════════════════
\newpage
\section*{Evidencia computacional — Capturas del programa}

Las siguientes capturas corresponden a los resultados generados por el
programa \texttt{integracion\_app.py}.

% ── Captura 1 ─────────────────────────────────────────────
\begin{figure}[H]
  \centering
  \includegraphics[width=0.82\linewidth]{01_datos.png}
  \caption{Datos de entrada — función, límites $a$, $b$ y valor de $N$.}
\end{figure}

% ── Captura 2 ─────────────────────────────────────────────
\begin{figure}[H]
  \centering
  \includegraphics[width=\linewidth]{02_tabla_trapecio.png}
  \caption{Tabla del Método Trapecial con coeficientes y sumatoria.}
\end{figure}

\newpage

% ── Captura 3 ─────────────────────────────────────────────
\begin{figure}[H]
  \centering
  \includegraphics[width=\linewidth]{03_tabla_newton_cotes.png}
  \caption{Tabla de Newton-Cotes: Simpson $\tfrac{1}{3}$, Regla $\tfrac{3}{8}$ y Boole.}
\end{figure}

% ── Captura 4 ─────────────────────────────────────────────
\begin{figure}[H]
  \centering
  \includegraphics[width=0.75\linewidth]{04_richardson.png}
  \caption{Aplicación del Límite Diferido de Richardson con $N$ y $2N$.}
\end{figure}

\newpage

% ── Captura 5 ─────────────────────────────────────────────
\begin{figure}[H]
  \centering
  \includegraphics[width=0.82\linewidth]{05_resultados.png}
  \caption{Resumen de resultados — comparación de los métodos aplicados.}
\end{figure}

% ── Captura 6 (extra, si aplica) ──────────────────────────
% \begin{figure}[H]
%   \centering
%   \includegraphics[width=0.82\linewidth]{06_extra.png}
%   \caption{Captura adicional.}
% \end{figure}

% ══════════════════  FIN  ═══════════════════════════════
\end{document}
